\documentclass[10pt]{article}
\usepackage[T1]{fontenc}
\usepackage[utf8]{inputenc}
\usepackage[french]{babel}
\usepackage{amsmath,amssymb,mathrsfs,mathpazo}
\usepackage{geometry}
\geometry{a4paper,margin=2cm}
\usepackage{graphicx,array,tabularx,tabularray,colortbl,multirow,multicol,booktabs,enumitem,pifont,ulem,fancybox,diagbox,cancel}
\usepackage{pstricks,pstricks-add,pst-plot,pst-tree,pst-node,pst-text,pst-func,pst-math,pst-bspline,pst-3dplot,pst-eucl,pst-solides3d}
\let\pscalculate\undefined
\usepackage{pst-bezier}
\usepackage{tikz,pgfplots,tkz-tab}
\usetikzlibrary{arrows,arrows.meta,patterns,calc,positioning,trees,decorations.markings,shapes.geometric,shapes.misc,shapes.arrows,intersections,decorations.pathmorphing,decorations.pathreplacing,backgrounds,shadows,babel}
\pgfplotsset{compat=1.18}
\usepgfplotslibrary{fillbetween}
\usepackage[np]{numprint}
\usepackage{listings,makecell,pas-tableur,twoopt}
\let\barre\undefined
\usepackage[pstricks]{tablvar}
\definecolor{darkgreen}{RGB}{0,100,0}
\providecommand{\red}{\color{red}}
\providecommand{\blue}{\color{blue}}
\providecommand{\green}{\color{green}}
\providecommand{\black}{\color{black}}
\providecommand{\white}{\color{white}}
\providecommand{\gray}{\color{gray}}
\providecommand{\cyan}{\color{cyan}}
\providecommand{\magenta}{\color{magenta}}
\providecommand{\yellow}{\color{yellow}}
\providecommand{\violet}{\color{violet}}
\providecommand{\orange}{\color{orange}}
\providecommand{\euro}{\texteuro}
\providecommand{\EUR}{\texteuro}
\providecommand{\ieme}{\textsuperscript{e}}
\providecommand{\degres}{\ensuremath{^\circ}}
\providecommand{\degree}{\ensuremath{^\circ}}
\providecommand{\vv}[1]{\overrightarrow{#1}}
\providecommand{\Coord}[1]{#1}
\providecommand{\decosix}{}
\providecommand{\starredbullet}{\ensuremath{\bullet}}
\pagestyle{empty}
\setlength{\parindent}{0pt}
\usepackage{adjustbox}\DeclareUnicodeCharacter{2236}{:}\DeclareUnicodeCharacter{2190}{\ensuremath{\leftarrow}}
\providecommand{\R}{\mathbb{R}}
\providecommand{\N}{\mathbb{N}}
\providecommand{\Z}{\mathbb{Z}}
\providecommand{\Q}{\mathbb{Q}}
\providecommand{\D}{\mathbb{D}}
\providecommand{\C}{\mathbb{C}}
\providecommand{\e}{{\mathrm{e}}}
\providecommand{\vect}[1]{\overrightarrow{#1}}
\providecommand{\vectt}[1]{\overrightarrow{\mathrm{#1}}}
\providecommand{\vvt}[1]{\overrightarrow{\mathrm{#1}}}
\providecommand{\np}[2][]{{#2~#1}}
\providecommand{\vv}[1]{{\overrightarrow{#1}}}
\renewcommand{\d}{{\mathrm{d}}}
\providecommand{\degres}{{^\circ}}
\providecommand{\degre}{{^\circ}}
\providecommand{\blue}{\color{blue}}
\providecommand{\red}{\color{red}}
\providecommand{\enskip}{\quad}
\providecommand{\llbracket}{\lbrack\!\lbrack}
\providecommand{\rrbracket}{\rbrack\!\rbrack}
\providecommand{\euro}{{\text{€}}}
\providecommand{\no}{{\text{n°}}}
\providecommand{\interval}[2]{{[#1;#2]}}
\providecommand{\intervFO}[2]{{[#1;#2[}}
\providecommand{\dotfill}{\ldots}

\providecommand{\ds}{\displaystyle}
\providecommand{\textding}[1]{#1}
\providecommand{\Oij}{\(\mathrm{O};\vec{\imath},\vec{\jmath}\)}
\providecommand{\Oijk}{\(\mathrm{O};\vec{\imath},\vec{\jmath},\vec{k}\)}
\providecommand{\Ouv}{\(\mathrm{O};\vec{u},\vec{v}\)}
\providecommand{\intervFF}[2]{\ensuremath{[#1;#2]}}
\renewcommand{\euro}{\ifmmode\text{\texteuro}\else\texteuro\fi}
\DeclareUnicodeCharacter{00B0}{\ifmmode^{\circ}\else\textdegree\fi}
\DeclareUnicodeCharacter{20AC}{\ifmmode\text{\texteuro}\else\texteuro\fi}
% Énoncés : Éducation nationale. Transcription et corrigés : APMEP, Denis Vergès et contributeurs. Découpage et classement : CoopMaths, Sébastien Lozano, Éric Elter et contributeurs. Adaptation de la présentation : Maths Explorer.
% CC BY-SA : https://coopmaths.fr/www/about/
% https://coopmaths.fr/alea/static/eam/2025/tex/eam_specifique_2025_06_sujet0v3_4_cor.tex
\begin{document}
\noindent\textit{APMEP / CoopMaths. CC BY-SA. Présentation adaptée par Maths Explorer.}\par\medskip
\medskip





\begin{enumerate}
\item Un employé reçoit des appels téléphoniques.

On estime que la probabilité qu'un appel dure plus de cinq minutes est égale à 0,3.\\
On suppose que les durées des différents appels sont indépendantes.\\
Ce matin, l'employé reçoit deux appels.

\smallskip 



\begin{list}{\textbullet}{On représente la situation au moyen d'un arbre de probabilité en appelant:}
\item $P$ l'événement \og l'appel a duré plus de 5 minutes \fg{};
\item $M$ l'événement \og l'appel n'a pas duré plus de 5 minutes \fg{}.
\end{list}

\begin{center}
\bigskip
  \pstree[treemode=R,nodesepA=0pt,nodesepB=4pt,levelsep=2.5cm]{\TR{}}
 {
 	\pstree[nodesepA=4pt]{\TR{$P$}\ncput*{$0,3$}}
 	  { 
 		  \TR{$P$}~{\blue $\longrightarrow PP\quad \red 0,3\times 0,3=0,09$}\ncput*{$0,3$}
 		  \TR{$M$}~{\blue $\longrightarrow PM\quad \red 0,3\times 0,7=0,21$}\ncput*{$0,7$}	   
 	  }
 	\pstree[nodesepA=4pt]{\TR{$M$}\ncput*{$0,7$}}
 	  {
 		  \TR{$P$}~{\blue $\longrightarrow MP\quad \red 0,7\times 0,3=0,21$}\ncput*{$0,3$}
 		  \TR{$M$}~{\blue $\longrightarrow MM\quad \red 0,7\times 0,7=0,49$}\ncput*{$0,7$}	   
      }
}
\bigskip
\end{center}

\textbf{Affirmation 1 :}

La probabilité que les deux appels durent tous les deux plus de cinq minutes est égale à 0,09.

\begin{tabular}{@{\hspace{0.03\linewidth}} | p{0.95\linewidth}}
L'événement \og les deux appels durent tous les deux plus de cinq minutes\fg{} correspond à l'évènement $PP$ dont la probabilité est: $0,3\times 0,3=0,09$.

\hfill\textbf{Affirmation 1 vraie}
\end{tabular}

\medskip 

\textbf{Affirmation 2 :}

La probabilité qu'un appel exactement sur les deux dure plus de cinq minutes est égale à 0,21.

\begin{tabular}{@{\hspace{0.03\linewidth}} | p{0.95\linewidth}}
L'événement \og un appel exactement sur les deux dure plus de cinq minutes\fg{} correspond à la réunion des évènements incompatibles $PM$ et $MP$.

La probabilité de chacun de ces événements est $0,3\times 0,7=0,21$, donc la probabilité cherchée est de $0,21+0,21=0,42$.

\hfill\textbf{Affirmation 2 fausse}
\end{tabular}


\medskip 

\item Le gérant d'une piscine s'intéresse à la présence de bactéries dans l'eau. Il effectue un prélèvement. Ce prélèvement montre que la concentration de bactéries est égale à $\np{1000}$ bactéries par millilitre. Le seuil maximal autorisé est égal à $\np{1500}$ bactéries par millilitre.

On admet que la concentration de bactéries est modélisée par la fonction $f$ définie sur l'intervalle~$[0~; ~+\infty[$ par :
$f(t) = 1,1^{t}$,
où $f(t)$ désigne la concentration, en milliers de bactéries par millilitre, et $t$ désigne la durée, en heure, écoulée depuis que le prélèvement a été effectué.

\smallskip 

\textbf{Affirmation 3 :}

La fonction $f$ est croissante sur l'intervalle $[0~;~+\infty[$.

\smallskip 

\begin{tabular}{@{\hspace{0.03\linewidth}} | p{0.95\linewidth}}
La fonction $f$ est une fonction exponentielle de type $t \longmapsto a^t$, avec $a=1,1$; or $1,1>1$ donc la fonction $f$ est strictement croissante.

\hfill\textbf{Affirmation 3 vraie}
\end{tabular}

\medskip 

\textbf{Affirmation 4 :}

La concentration de bactéries deux heures après le prélèvement est inférieure au seuil maximal autorisé.

\smallskip

\begin{tabular}{@{\hspace{0.03\linewidth}} | p{0.95\linewidth}}
$f(2)=1,1^2=1,21$; donc la concentration de bactéries deux heures après le prélèvement est de $\np{1210}$ bactéries, donc inférieure au seuil maximal autorisé. 

\hfill\textbf{Affirmation 4 vraie}
\end{tabular}

\end{enumerate}

\end{document}
